% ============================================================================
% Elementos pré-textuais: capa, folha de rosto, resumo, abstract e sumário
% ============================================================================

% ------------------------------------------------------------------- CAPA ----
\begin{titlepage}
  \begin{center}
    \fontsize{12}{18}\selectfont
    \textbf{UNIVERSIDADE PAULISTA -- UNIP}\\[0.2\baselineskip]
    Projeto Integrado Multidisciplinar\\[0.2\baselineskip]
    Curso Superior de Tecnologia em\\ Análise e Desenvolvimento de Sistemas

    \vspace{2.2cm}

    \textbf{[NOME DO ALUNO] -- [RA]}\\[0.25\baselineskip]
    \textbf{[NOME DO ALUNO] -- [RA]}

    \vspace{2.4cm}

    {\fontsize{14}{21}\selectfont\bfseries
      AGROCONECTA: PLATAFORMA CORPORATIVA INTEGRADA PARA GESTÃO DA CADEIA
      DE SUPRIMENTOS E INCLUSÃO DIGITAL NA AGRICULTURA FAMILIAR\par}

    \vfill

    São Paulo -- SP\\
    2026
  \end{center}
\end{titlepage}

% --------------------------------------------------------- FOLHA DE ROSTO ---
\begin{titlepage}
  \begin{center}
    \fontsize{12}{18}\selectfont
    \textbf{[NOME DO ALUNO] -- [RA]}\\[0.25\baselineskip]
    \textbf{[NOME DO ALUNO] -- [RA]}

    \vspace{2.6cm}

    {\fontsize{14}{21}\selectfont\bfseries
      AGROCONECTA: PLATAFORMA CORPORATIVA INTEGRADA PARA GESTÃO DA CADEIA
      DE SUPRIMENTOS E INCLUSÃO DIGITAL NA AGRICULTURA FAMILIAR\par}

    \vspace{2.5cm}
  \end{center}

  \noindent\hspace*{5.5cm}%
  \begin{minipage}{9.5cm}
    \raggedright
    \fontsize{11}{16}\selectfont
    Projeto Integrado Multidisciplinar apresentado à Universidade Paulista --
    UNIP, como requisito parcial para a obtenção do título de tecnólogo em
    Análise e Desenvolvimento de Sistemas.

    \vspace{0.6\baselineskip}
    Orientador(a): [NOME DO ORIENTADOR]
  \end{minipage}

  \vfill
  \begin{center}
    \fontsize{12}{18}\selectfont
    São Paulo -- SP\\
    2026
  \end{center}
\end{titlepage}

% ----------------------------------------------------------------- RESUMO ----
\setcounter{page}{3}
\pagestyle{vazio}
\begingroup
\setstretch{1}
\titulopretextual{RESUMO}

A transformação digital de cadeias produtivas agroalimentares exige soluções que
conectem planejamento de negócio, aplicações corporativas, banco de dados,
infraestrutura e gestão ágil em um único ecossistema. Este trabalho apresenta o
AgroConecta, plataforma corporativa integrada para a gestão da cadeia de
suprimentos e a inclusão digital da agricultura familiar. O problema identificado
é a assimetria de informação entre cooperativas, produtores familiares,
compradores institucionais e transportadores, que provoca perdas de produção,
fragilidade na rastreabilidade e baixa capacidade de escoamento. A metodologia
adotou pesquisa aplicada de natureza qualitativa, apoiada em levantamento
documental, modelagem de negócio e desenvolvimento incremental com práticas ágeis
(Scrum e Kanban). A solução foi construída sobre uma arquitetura em camadas com
.NET 10 e ASP.NET Core, oferecendo aplicação Web com autenticação por cookie e
API REST protegida por JWT, além de um aplicativo mobile em React Native e Expo
com operação offline. O banco de dados relacional em SQL Server 2022 contempla
quatorze tabelas de negócio, três procedimentos armazenados e três gatilhos de
auditoria e movimentação de estoque, validados por scripts completos. Os
resultados incluem um fluxo ponta a ponta de produto, lote, pedido, rota, entrega
com evidência fotográfica e geolocalização, e faturamento, com 159 testes
automatizados aprovados e cobertura superior a 90\% nas camadas de domínio e
aplicação. Conclui-se que a plataforma atende aos requisitos técnicos e
organizacionais do projeto, amplia a transparência da cadeia e cria condições
concretas de inclusão produtiva para agricultores familiares.

\vspace{0.5\baselineskip}
\noindent\textbf{Palavras-chave:} Agricultura familiar. Transformação digital.
Plataforma corporativa. Desenvolvimento Web. Aplicativo mobile. Banco de dados.
\endgroup

% ---------------------------------------------------------------- ABSTRACT ---
\newpage
\begingroup
\setstretch{1}
\begin{otherlanguage}{english}
\titulopretextual{ABSTRACT}

The digital transformation of agrifood supply chains requires solutions that
connect business planning, corporate applications, databases, infrastructure and
agile management in a single ecosystem. This paper presents AgroConecta, an
integrated corporate platform for supply chain management and digital inclusion
of family farming. The identified problem is the information asymmetry among
cooperatives, family producers, institutional buyers and transporters, which
causes production losses, weak traceability and low market access. The
methodology adopted applied qualitative research supported by document analysis,
business modeling and incremental development with agile practices (Scrum and
Kanban). The solution was built on a layered architecture with .NET 10 and
ASP.NET Core, providing a Web application with cookie authentication and a REST
API protected by JWT, as well as a React Native and Expo mobile application with
offline operation. The relational database on SQL Server 2022 comprises fourteen
business tables, three stored procedures and three triggers for auditing and
inventory movement, validated by complete scripts. The results include an
end-to-end flow of product, harvest lot, order, route, delivery with photographic
and geolocation evidence, and invoicing, with 159 automated tests passing and
coverage above 90\% in the domain and application layers. It is concluded that
the platform meets the technical and organizational requirements of the project,
increases supply chain transparency and creates concrete conditions for
productive inclusion of family farmers.

\vspace{0.5\baselineskip}
\noindent\textbf{Keywords:} Family farming. Digital transformation. Corporate
platform. Web development. Mobile application. Database.
\end{otherlanguage}
\endgroup

% ---------------------------------------------------------------- SUMÁRIO ----
\newpage
\begingroup
\setstretch{1}
\tableofcontents
\endgroup
